% kl_theory.tex
% Theory section for KL-divergence rank-preserving calibration
% To be included in the follow-up paper

%----------------------------------------------------------------------
% BREGMAN KKT UMBRELLA THEOREM
%----------------------------------------------------------------------

\subsection{Bregman Umbrella Framework}

We present a unified framework for rank-preserving calibration using Bregman divergences,
which includes both Euclidean (squared loss) and KL (relative entropy) as special cases.

Let $D_\phi(Q,R) = \sum_{i,j} \left[\phi(Q_{ij}) - \phi(R_{ij}) - \nabla\phi(R_{ij})(Q_{ij} - R_{ij})\right]$
denote the Bregman divergence generated by a strictly convex function $\phi$.

\begin{theorem}[Bregman KKT Umbrella]
\label{thm:bregman-kkt}
Consider the penalized rank-preserving calibration problem:
\begin{equation}
\min_{Q \in \mathcal{C}(T)} D_\phi(Q, R) + \lambda V_A(Q)
\label{eq:bregman-problem}
\end{equation}
where $\mathcal{C}(T) = \{Q : Q \mathbf{1} = \mathbf{1}, \, Q^\top \mathbf{1} = T, \, Q \geq 0\}$
is the transportation polytope with row-simplex and column-sum constraints,
and $V_A(Q) = \sum_{j=1}^J \sum_{r=1}^{N-1} \omega_{jr} [(D S_j Q_{\cdot j})_r]_+$
is the weighted rank violation penalty with $S_j$ permuting column $j$ by anchor $A$.

The KKT conditions for \eqref{eq:bregman-problem} are:
\begin{equation}
\nabla \phi(Q^\star_{ij}) = \nabla \phi(R_{ij}) - \nu_i - \mu_j - [S_j^\top D^\top \alpha_j]_i
\label{eq:bregman-kkt}
\end{equation}
where $\nu_i$ and $\mu_j$ are dual variables for row and column constraints respectively,
and $0 \leq \alpha_{jr} \leq \lambda \omega_{jr}$ are the rank violation multipliers.

\medskip
\noindent\textbf{Special Cases:}

\begin{enumerate}[label=(\alph*)]
\item \textbf{Euclidean} ($\phi(x) = \frac{1}{2}x^2$, so $\nabla\phi(x) = x$):
\begin{equation}
Q^\star_{ij} = R_{ij} - \nu_i - \mu_j - [S_j^\top D^\top \alpha_j]_i
\end{equation}
The solution is \emph{additive} in dual potentials.

\item \textbf{KL} ($\phi(x) = x \log x - x$, so $\nabla\phi(x) = \log x$):
\begin{equation}
Q^\star_{ij} = R_{ij} \exp\left(-\nu_i - \mu_j - [S_j^\top D^\top \alpha_j]_i\right)
\end{equation}
The solution is \emph{multiplicative} in dual potentials.
\end{enumerate}
\end{theorem}

\begin{proof}
The Lagrangian for \eqref{eq:bregman-problem} is:
\begin{align}
\mathcal{L}(Q, \nu, \mu, \alpha) &= D_\phi(Q, R) + \lambda V_A(Q) \\
&\quad + \sum_i \nu_i \left(\sum_j Q_{ij} - 1\right) + \sum_j \mu_j \left(\sum_i Q_{ij} - T_j\right) \nonumber
\end{align}

Taking the derivative with respect to $Q_{ij}$ and using the subgradient of $[\cdot]_+$:
\[
\frac{\partial \mathcal{L}}{\partial Q_{ij}} = \nabla\phi(Q_{ij}) - \nabla\phi(R_{ij}) + \nu_i + \mu_j + [S_j^\top D^\top \alpha_j]_i = 0
\]
Rearranging yields \eqref{eq:bregman-kkt}. The special cases follow by substituting
$\nabla\phi(x) = x$ (Euclidean) and $\nabla\phi(x) = \log x$ (KL).
\end{proof}

%----------------------------------------------------------------------
% KL COLUMN PROJECTION LEMMA
%----------------------------------------------------------------------

\subsection{KL Isotonic Projection with Sum Constraint}

The column subproblem for KL calibration requires projecting onto the intersection
of the isotonic cone and a hyperplane (sum constraint). Unlike the Euclidean case
which uses arithmetic mean pooling and additive shifts, KL requires geometric mean
pooling and multiplicative scaling.

\begin{lemma}[KL Column Projection]
\label{lem:kl-column}
For KL divergence, the column subproblem
\begin{equation}
\min_{x_1 \leq \cdots \leq x_N, \, \sum_i x_i = T} \sum_{i=1}^N \left[x_i \log\frac{x_i}{w_i} - x_i + w_i\right]
\label{eq:kl-column}
\end{equation}
has solution obtained by:
\begin{enumerate}
\item Apply generalized PAV with \textbf{geometric mean} pooling
\item Apply \textbf{multiplicative rescaling}: $x^\star = y \cdot (T / \sum_i y_i)$
\end{enumerate}

\noindent\textit{Contrast with Euclidean:} PAV (arithmetic mean) + additive shift.
\end{lemma}

\begin{proof}
\textbf{Step 1: Geometric Mean Pooling.}

Consider a block $B$ of consecutive indices that must be pooled to satisfy isotonicity.
The KL-optimal constant value $z_B$ for this block minimizes:
\[
\sum_{i \in B} \left[z_B \log\frac{z_B}{w_i} - z_B + w_i\right]
\]
Taking the derivative with respect to $z_B$:
\[
\frac{\partial}{\partial z_B} = \sum_{i \in B} \left[\log z_B - \log w_i + 1 - 1\right] = |B| \log z_B - \sum_{i \in B} \log w_i = 0
\]
Solving:
\[
\log z_B = \frac{1}{|B|} \sum_{i \in B} \log w_i \quad\Rightarrow\quad z_B = \exp\left(\frac{1}{|B|} \sum_{i \in B} \log w_i\right) = \left(\prod_{i \in B} w_i\right)^{1/|B|}
\]
This is the \textbf{geometric mean}, contrasting with the arithmetic mean for Euclidean.

\medskip
\textbf{Step 2: Multiplicative Rescaling.}

After PAV, let $y = (y_1, \ldots, y_N)$ be the isotonic solution. To satisfy
$\sum_i x_i = T$, we seek a scalar $c$ such that $\sum_i c \cdot y_i = T$.
Thus $c = T / \sum_i y_i$ and $x^\star_i = c \cdot y_i$.

This multiplicative scaling preserves isotonicity since the isotone cone
$\{x : x_1 \leq \cdots \leq x_N\}$ is closed under positive scalar multiplication.

For Euclidean, the optimal adjustment is additive: $x^\star_i = y_i + (T - \sum_i y_i)/N$.
\end{proof}

%----------------------------------------------------------------------
% ANCHOR-REFERENCE DECOUPLING
%----------------------------------------------------------------------

\subsection{Anchor-Reference Decoupling}

A key innovation in our KL framework is the separation of two roles:
\begin{itemize}
\item \textbf{Reference $R$}: The distribution from which we measure KL divergence.
\item \textbf{Anchor $A$}: The distribution that determines rank orderings to preserve.
\end{itemize}

\begin{definition}[Decoupled KL Calibration]
Given input $P$, reference $R$, anchor $A$, and target marginals $T$:
\begin{equation}
Q^\star = \arg\min_{Q \in \mathcal{C}(T)} \text{KL}(Q \| R) \quad\text{s.t.}\quad Q_{\cdot j} \text{ is isotonic w.r.t. } A_{\cdot j} \text{ ordering}
\end{equation}
\end{definition}

When $R = A = P$, this reduces to standard rank-preserving calibration.
Decoupling enables scenarios where:
\begin{itemize}
\item We want to measure divergence from one model ($R$) while preserving
      rankings from another ($A$) --- e.g., ensemble calibration.
\item The reference $R$ is a prior or historical distribution, while $A$
      comes from the current classifier.
\item Label shift adjustment where $R$ reflects source domain and $A$
      reflects target domain ordering preferences.
\end{itemize}

%----------------------------------------------------------------------
% PARETO FRONTIER
%----------------------------------------------------------------------

\subsection{Pareto Frontier Characterization}

Rather than reporting a single solution for a tuned $\lambda$, we advocate
computing the full Pareto frontier of KL divergence versus rank violation.

\begin{proposition}[Pareto Monotonicity]
For the soft problem \eqref{eq:bregman-problem}, let $Q^\star(\lambda)$ denote
the solution for penalty parameter $\lambda$. Then:
\begin{enumerate}
\item $V_A(Q^\star(\lambda))$ is non-increasing in $\lambda$.
\item $\text{KL}(Q^\star(\lambda) \| R)$ is non-decreasing in $\lambda$.
\end{enumerate}
Thus $\{(\text{KL}(Q^\star(\lambda)\|R), V_A(Q^\star(\lambda))) : \lambda \geq 0\}$
traces the Pareto frontier.
\end{proposition}

\begin{proof}
Standard results from multi-objective optimization: increasing $\lambda$
places more weight on rank violation, driving $V_A$ down at the cost of
increased KL divergence from the reference.
\end{proof}

The Pareto frontier allows practitioners to choose an operating point based on
their relative preference for calibration accuracy versus rank preservation,
rather than relying on a single ``optimal'' $\lambda$.
